% !PRINTED 570
\weberpage{570}
%<<CONT
jede beliebig gegebene natürliche Primzahl $p$ in ihre Primfactoren
wirklich zu zerlegen, und damit alle Primideale des Körpers $\Omega$,
oder wenigstens Repräsentanten aller Primideale wirklich
darzustellen\footnote{Wir folgen hier einer Arbeit von \gesperrt{Hensel}: \glqq Untersuchung der
Fundamentalgleichung einer Gattung für eine reelle Primzahl als Modul
und Bestimmung der Theiler der Discriminante\grqq{} (\gesperrt{Crelle}'s Journ., Bd. 113).
Zu erwähnen sind auch die anderen Arbeiten von \gesperrt{Hensel}, Ebend.,
Bd. 101, 104, 105, 113.}.

Es sei $\mathfrak p$ ein beliebiges Primideal vom Grade $f$, so dass
\weq{1}{N(\mathfrak p) = p^f}
ist, worin $p$ eine natürliche Primzahl bedeutet, die durch den
Primfactor $\mathfrak p$ theilbar ist. $f$ ist ein positiver Exponent
$\leqq n$. Die kleinste ganze rationale Zahl, die durch $\mathfrak p$
theilbar ist, ist $p$, und jede andere ganze rationale Zahl, die durch
$\mathfrak p$ theilbar ist, ist daher auch durch $p$ theilbar
(\S.\,138).

Wir müssen nun auch Congruenzen mit dem Modul $\mathfrak p$ zwischen
ganzen Functionen beliebiger Variablen mit Coëfficienten in
$\mathfrak o$ betrachten und bemerken, dass zwei solche Functionen
nach \S.\,143, 2. \gesperrt{dann und nur dann congruent sind, wenn die
entsprechenden Coëfficienten congruent sind.}

Den Körper der rationalen Zahlen bezeichnen wir mit $R$ und nennen
demnach eine Function mit rationalen Coëfficienten auch eine \gesperrt{Function
in} $R$.

Die Basisform von $\mathfrak o$
\weq{2}{\tau = \omega_1 t_1 + \omega_2 t_2 + \cdots + \omega_n t_n}
genügt, wie wir im \S.\,146 gesehen haben, einer Gleichung
$n^{\text{ten}}$ Grades $F(\tau) = 0$, deren Coëfficienten ganze
rationale Functionen von $t_1, t_2, \ldots, t_n$ sind. Es ist also
$F(\tau)$ jedenfalls durch $\mathfrak p$ theilbar, und daraus folgt,
dass es ganze Functionen $\Phi(t)$ in $R$ giebt, die ausser $t$
irgend welche Variable enthalten können, die durch die Substitution
$t = \tau$ in durch $\mathfrak p$ theilbare Functionale in
$\mathfrak o$ übergehen, die also der Congruenz
\weq{3}{\Phi(\tau) \equiv 0\ (\mathrm{mod}\ \mathfrak p)}
genügen, und wir werden also sagen können, $\tau$ \gesperrt{ist eine Wurzel der
Congruenz}
\weq{4}{\Phi(t) \equiv 0\ (\mathrm{mod}\ \mathfrak p).}

Die Function $\Phi$ wird gewiss die Variablen $t_1, t_2, \ldots, t_n$
enthalten müssen; sie kann aber auch noch andere Variable ent-
%CONT>>
